\clearpage
\fancyhead[L]{\small EFFEMERIDI ASTRONOMICHE DI MILANO PER L'ANNO 1817.}
\hypertarget{native-606-lower}{}
\begin{center}\small Göttingische gelehrte Anzeigen. Stück 144. Seite 1433..1438. 1817 September 8.\\[3pt]\rule{.72\linewidth}{.3pt}\end{center}

{\setlength{\rightskip}{0pt plus 3em}\emph{Effemeridi astronomiche di Milano per l'anno \textup{1817} calcolate da} \textsc{Francesco Carlini} \emph{ed} \textsc{Enrico Brambilla}. \emph{Con Appendice.} Aus der K. K. Druckerei. Mailand 1816. 108 und 116 Seiten in Octav.\par}

Da die Einrichtung des Kalenders in diesem Jahrgange, und die beigefügten Zusätze ganz dieselben sind, wie bei dem vorhergehenden, so haben wir diesmal blos die Aufsätze des Anhanges anzuzeigen. In dem ersten liefert uns Herr \textsc{Oriani} das Verzeichniss der Declinationen von 40 Sternen aus den in frühern Jahrgängen der Ephemeriden abgedruckten Beobachtungen mit dem dreifüssigen \textsc{Reichenbach}'schen Vervielfältigungskreise. Die beigefügte Vergleichung mit den Angaben von \textsc{Maskelyne}, \textsc{Piazzi} und \textsc{Pond} ist wegen der von diesen Astronomen angewandten Instrumente merkwürdig; in Rücksicht auf die erreichbare Genauigkeit setzt Hr. \textsc{Oriani} den dreifüssigen \textsc{Reichenbach}'schen Wiederholungskreis dem fünffüssigen \textsc{Troughton}'schen Mauerkreis in Greenwich ungefähr gleich, und in Rücksicht auf die Bequemlichkeit des Beobachtens räumt er gleichfalls letzterm Instrumente keinen Vorzug ein, indem das, was ein viermaliges Wiederholen an Zeit mehr kostet, durch das am Wiederholungskreise bequemere Ablesen ersetzt werde. Hierbei ist freilich der zur Berechnung der ausser dem Meridian gemachten Beobachtungen erforderliche Zeitaufwand nicht in Anschlag gebracht. --- Der zweite Aufsatz, von \textsc{Ottaviano Fabrizio Mossotti}, neue Analyse der Aufgabe, die Bahnen der Himmelskörper zu bestimmen, verdient wegen der Wichtigkeit des Gegenstandes eine etwas umständliche Erwähnung. Die ganze Behandlung ist rein analytisch und mit Geschicklichkeit durchgeführt: inzwischen wird durch die grosse Menge der Zeichen die Festhaltung des Fadens, die Uebersicht des Ganzen und die klare Auffassung dessen, worauf es eigentlich ankommt, etwas erschwert. Das Wesentliche beruht auf folgendem: Wenn die halben Parameter der zu bestimmenden Bahn und der Erdbahn mit $p$ und $P$; die Producte von $\BareRoot{p}$ \SourceLeaf{607}in die Cosinus der Neigung der Ebene der unbekannten Bahn gegen drei feste willkührliche Ebnen mit
$c$, $c'$, $c''$, und eben so die Producte von $\BareRoot{P}$ in die Cosinus der Neigungen der Erdbahn gegen dieselben
Fundamentalebenen mit $C$, $C'$, $C''$ bezeichnet werden, so findet Herr \textsc{Mossotti} drei linearische Gleichungen von der Form
\[
A(c-C)+A'(c'-C')+A''(c''-C'') = 0
\] 
wo die Coëfficienten bekannte Grössen sind. Diese Gleichungen sind jedoch nur näherungsweise richtig,
in so fern die Zwischenzeiten nicht unendlich klein sind, und gründen sich eigentlich blos darauf, dass
die drei Flächensegmente, welche entstehen, wenn die drei Oerter paarweise durch Sehnen verbunden
werden, sowohl bei der unbekannten Bahn, als bei der Erdbahn sehr nahe den Würfeln der Zwischenzeiten proportional sind. Eben deswegen, weil jene Gleichungen nur genähert wahr sind, ist es nicht
verstattet, sie so zu combiniren, dass eine der Grössen $c-C$, $c'-C'$, $c''-C''$, eliminirt würde, sondern
es lässt sich zeigen, dass alle drei Gleichungen eigentlich nur für Eine gelten dürfen, und Herr \textsc{Massotti} behält daher auch nur Eine von ihnen bei. Da nun diese für sich allein nicht weiter führen kann,
so nimmt Herr \textsc{Massotti} eine vierte Beobachtung zu Hülfe, die, mit zweien der vorigen verbunden, eine
ganz ähnliche Gleichung liefert, deren Verbindung mit der aus den drei ersten Beobachtungen gefolgerten allerdings rechtmässig ist, und zur Bestimmung des Verhältnisses der drei Grössen $c-C$, $c'-C'$,
$c''-C''$ unter einander benutzt wird, so dass die Resultate in dieser Gestalt erscheinen
\[
c'-C' = M(c-C), \quad c''-C'' = N(c-C)
\]
Durch Benutzung des oben ausgesprochenen Princips findet hernach Herr \textsc{Massotti} auch das Verhältniss
der drei Abstände des Himmelskörpers von der Erde zu einer der drei Grössen $c-C$, $c'-C'$, $c''-C''$,
und endlich die absoluten Werthe dieser drei Grössen selbst. Da nun $C$, $C'$, $C''$ für sich bekannt sind,
so ergeben sich die Werthe von $c$, $c'$, $c''$, und damit $p = cc+c'c'+c''c''$, so wie aus
\[
\frac{c}{\BareRoot{p}}, \quad \frac{c'}{\BareRoot{p}}, \quad \frac{c''}{\BareRoot{p}}
\]
die Grössen, von welchen die Lage der Bahn des Himmelskörpers abhängt. Die weitere Bestimmung
der übrigen Elemente beruhet auf bekannten Sätzen und hat nichts Eigenthümliches.

Sollen wir nun unser Urtheil über diese neue Methode hier abgeben, so erkennen wir zuvörderst
mit Vergnügen an, dass die erwähnte linearische Gleichung zwischen $c-C$, $c'-C'$, $c''-C''$ merkwürdig, und dass es interessant ist, die Möglichkeit einer durchgehends blos linearischen Bestimmung einer unbekannten Bahn entwickelt zu sehen. Eine andere Frage ist nun aber, ob eine solche Methode
zur wirklichen practischen Anwendung zu empfehlen sei. Es ist klar, dass hier nicht alles benutzt ist,
und nicht benutzt werden sollte, was in der Theorie der Bewegung in Kegelschnitten liegt, und zwar
wird namentlich, um es im Geiste einer Näherungsmethode auszusprechen, die Bedingung ganz ignorirt, dass die Quotienten, wenn jene Segmente mit den Würfeln der Zwischenzeiten dividirt werden,
sehr nahe gleich sein müssen der Grösse $\dfrac{k^3}{12r^3}$, wo $k$ in der Bedeutung der Theoria Motus Corporum
Coelestium zu nehmen ist, und $r$ den Abstand des Himmelskörpers von der Sonne bedeutet. Daher
also die Nothwendigkeit, eine vierte Beobachtung zu Hülfe zu nehmen bei einer Aufgabe, wo drei
Beobachtungen schon die vollständige Auflösung einschliessen. Auch abgesehen davon, dass dies in
theoretischer Rücksicht nicht geziemend ist, wird dies Verfahren in practischer Hinsicht deswegen sehr
misslich, weil eigentlich hier die Bestimmung der Bahn auf Grössen der \emph{dritten} Ordnung beruhet, und
daher die unvermeidlichen Beobachtungsfehler einen unverhältnissmässigen Einfluss auf die Resultate
äussern müssen, der desto grösser sein wird, weil die Methode, als blosse Näherungsmethode, nur kurze \SourceLeaf{608}Zwischenzeiten anzuwenden gestatten kann. Diese Methode wird daher im Allgemeinen nur wenig zuverlässige Resultate geben können, wo die vollständige Benutzung dreier Beobachtungen schon zu einer sehr genäherten Bestimmung führen kann.

Endlich müssen wir noch bemerken, dass, wenn man einmal blos von jenem Princip der Proportionalität der Segmente zu den Würfeln der Zwischenzeiten ausgehen will, man denselben Erfolg einer blos linearischen Bestimmung der Bahn aus vier Beobachtungen auf eine einfachere Art haben kann. Jenes Princip gibt nämlich aus drei Beobachtungen mit leichter Mühe eine gleichfalls linearische Gleichung zwischen zwei Abständen von der Erde, und zwischen denselben Abständen findet man auf dieselbe Weise eine zweite ähnliche Gleichung, indem man die zugehörigen beiden Beobachtungen mit einer vierten Beobachtung verbindet: die Elimination gibt dann diese Abstände selbst, eben so genau wie die Werthe, welche dieselben bei der Anwendung von \textsc{Mossotti}'s Methode erhalten, oder vielmehr eigentlich dieselben Werthe. Sobald diese Abstände bekannt sind, hat die Bestimmung der ganzen Bahn bekanntlich keine Schwierigkeit. Die weitere Entwickelung dieses Gegenstandes, wozu natürlich hier nicht der Ort ist, müssen wir auf eine andere Gelegenheit versparen. Nur das Eine müssen wir noch bemerken, dass der im vorliegenden Jahrgange der Ephemeriden abgedruckte 80 Seiten starke Aufsatz nur das Theoretische von \textsc{Mossotti}'s Methode enthält, und also eine versprochene Anwendung auf den \textsc{Halley}'schen Cometen wahrscheinlich im nächsten Jahrgange nachfolgen wird. --- Den Schluss des vorliegenden Jahrganges macht ein kleiner Artikel von \textsc{Carlini}, worin einige Fehler in \textsc{Delambre}'s Tafeln für die Jupiterssatelliten berichtigt werden.
